Con \(t\) expresado en segundos y \(x\) en metros,
\[
x(t)=
\begin{cases}
\dfrac14+\dfrac7{20}\cos(2t),&0\le t\le\pi/2,\\[2pt]
-\dfrac1{10},&t\ge\pi/2.
\end{cases}
\]
La primera parada posterior a la liberación ocurre en \(t=\pi/2\,\mathrm s\). La fuerza elástica es entonces \(0.40\,\mathrm N\) hacia la derecha; el rozamiento estático puede equilibrarla y la corredera permanece detenida. La energía disminuye de \(0.72\) a \(0.02\,\mathrm J\), con trabajo de rozamiento \(-0.70\,\mathrm J\).
